\documentclass[11pt]{article}
\usepackage[utf8]{inputenc}
\usepackage[T1]{fontenc}
\usepackage{amsmath,amssymb}
\usepackage{graphicx}
\usepackage{booktabs}
\usepackage{hyperref}
\usepackage[margin=1in]{geometry}
\usepackage{natbib}
\bibliographystyle{unsrtnat}

\title{A SUMO E3 ligase promotes long non-coding RNA transcription to regulate small RNA-directed DNA elimination}

\author{
Salman Shehzada$^{1}$ \and Kazufumi Mochizuki$^{1,*}$\\
\\
\small $^{1}$Institute of Human Genetics, University of Montpellier, CNRS, Montpellier, France\\
\small $^{*}$Corresponding author
}

\date{}

\begin{document}
\maketitle

\begin{abstract}
Small RNAs target their complementary chromatin regions for gene silencing through nascent long non-coding RNAs (lncRNAs). In the ciliated protozoan \textit{Tetrahymena}, the interaction between Piwi-associated small RNAs (scnRNAs) and the nascent lncRNA transcripts from the somatic genome has been proposed to induce target-directed small RNA degradation (TDSD), and scnRNAs not targeted for TDSD later target the germline-limited sequences for programmed DNA elimination. In this study, we show that the SUMO E3 ligase Ema2 is required for the accumulation of lncRNAs from the somatic genome and thus for TDSD and completing DNA elimination to make viable sexual progeny. Ema2 interacts with the SUMO E2 conjugating enzyme Ubc9 and enhances SUMOylation of the transcription regulator Spt6. We further show that Ema2 promotes the association of Spt6 and RNA polymerase II with chromatin. These results suggest that Ema2-directed SUMOylation actively promotes lncRNA transcription, which is a prerequisite for communication between the genome and small RNAs.
\end{abstract}

\section{Introduction}

Small RNAs of approximately 20--30 nucleotides that are complexed with Argonaute family proteins target either mRNAs for post-transcriptional silencing or chromatin regions for transcriptional and co-transcriptional silencing \citep{holoch2015}. For the latter process, small RNAs are generally considered to recognize their genomic targets via nascent long non-coding RNAs (lncRNAs) to induce heterochromatin formation \citep{johnson2002,kloc2008}. Therefore, chromatin regions targeted for silencing by small RNAs must be paradoxically transcribed to provide nascent lncRNAs.

In fission yeast, small interfering RNAs (siRNAs) mediate the deposition of histone H3 lysine 9 methylation (H3K9me) to establish heterochromatin. While the HP1 protein Swi6 binds to H3K9me2/3 for transcriptional silencing, phosphorylation of histone H3 serine 10 at the M phase of the cell cycle evicts Swi6, passively allowing lncRNA transcription \citep{chen2008,kloc2008}. In contrast, the HP1 paralog Rhino in fruit flies is specifically enriched at PIWI-interacting RNA (piRNA) clusters by binding to H3K9me2/3 \citep{leThomas2014} and actively recruits dedicated transcription machinery \citep{andersen2017}.

The ciliated protozoan \textit{Tetrahymena thermophila} possesses two types of nuclei: the transcriptionally active somatic macronucleus (MAC) and the germline micronucleus (MIC). During sexual reproduction (conjugation), the old MAC is degraded and a new MAC develops from a copy of the newly formed MIC. In this process, internal eliminated sequences (IESs) are removed from the developing new MAC genome through programmed DNA elimination \citep{austerberry1984}.

This process is guided by a class of Piwi-associated small RNAs called scnRNAs. scnRNAs are produced from the MIC genome, including both IESs and their surrounding macronucleus-destined sequences (MDSs) \citep{mochizuki2002,schoeberl2012}. In the parental MAC, scnRNAs complementary to the somatic (MDS) genome interact with nascent lncRNAs (pMAC-lncRNAs), and this interaction induces target-directed small RNA degradation (TDSD) \citep{aronica2008,noto2018}. scnRNAs not targeted for TDSD selectively accumulate and later direct the elimination of IESs from the developing new MAC.

The production of pMAC-lncRNAs in the parental MAC is essential for this genome-scanning mechanism, yet the molecular mechanisms that promote this transcription remain poorly understood. Here, we identify the SUMO E3 ligase Ema2 as a factor that actively promotes pMAC-lncRNA transcription through SUMOylation of the conserved transcription regulator Spt6, thereby facilitating the chromatin association of Spt6 and RNA polymerase II (RNAPII).

\section{Results}

\subsection{Ema2 is exclusively expressed during conjugation and localized in the MAC}

As part of a systematic investigation into genes highly upregulated during conjugation \citep{loidl2021}, we explored the function of \textit{EMA2} (TTHERM\_00113330). \textit{EMA2} mRNA is exclusively expressed during conjugation (Figure~1A). The Ema2 protein tagged with HA at the endogenous \textit{EMA2} locus was not detectable in vegetative cells and first appeared in the MAC during conjugation at 3 hours post-induction of mating (hpm) and 6~hpm (Figure~1B). At the onset of new MAC development, Ema2 disappeared from the parental MAC and appeared in the new MAC at 8~hpm, before fading away at later stages. This conjugation-specific expression and localization switch from the parental to the new MAC are reminiscent of factors involved in DNA elimination, such as the Piwi protein Twi1 and PRC2 \citep{mochizuki2002,liu2007,noto2010}.

\subsection{Ema2 is required for completing DNA elimination}

DNA elimination in exconjugants at 36~hpm was analyzed by DNA fluorescent in situ hybridization (FISH) using probes complementary to the transposable element Tlr1 \citep{wuitschick2002}. DNA elimination is completed by approximately 14--18~hpm in wild-type cells \citep{austerberry1984,mutazono2019}, and the Tlr1 element was detected only in the MICs in the exconjugants from wild-type cells (Figure~2A). In contrast, the Tlr1 element was detected in both the MICs and the MACs in the exconjugants from the \textit{EMA2} somatic KO strains (Figure~2A). The intensity of the FISH signal in the new MACs was lower in the exconjugants from the \textit{EMA2} KO cells than in those from the \textit{TWI1} KO cells, in which DNA elimination is known to be completely blocked \citep{noto2015}. Therefore, DNA elimination was partially blocked in the exconjugants of \textit{EMA2} KO cells. Consistent with the requirement of DNA elimination for the viability of sexual progeny \citep{cheng2010,vogt2013}, \textit{EMA2} KO cells did not produce viable progeny (Figure~2C).

\subsection{Ema2 is required for target-directed small RNA degradation (TDSD)}

We examined whether Ema2 is involved in TDSD. Northern blot analysis using the Mi-9 probe, complementary to a repetitive sequence found in both IESs and MDSs \citep{aronica2008}, showed that in wild-type cells, Mi-9-complementary scnRNAs were detected at 3~hpm, reduced at 4.5~hpm, and became undetectable at 6~hpm due to TDSD (Figure~3A). In contrast, these scnRNAs remained detectable at 6~hpm in \textit{EMA2} KO cells, indicating defective TDSD.

Small RNA sequencing confirmed this result at the genome-wide level. While scnRNAs mapping to MDS tiles were greatly reduced from 3 to 8~hpm in wild-type cells due to TDSD, those in the \textit{EMA2} KO cells were only slightly reduced (Figure~3B). This TDSD defect appeared milder than in \textit{EMA1} KO cells, where MDS-complementary scnRNAs remained constant until 8~hpm, indicating that Ema1 and Ema2 act differently in TDSD. In contrast, scnRNAs complementary to IES tiles were comparable across all strains at both time points, consistent with the selective targeting of MDS-complementary scnRNAs for TDSD.

\subsection{Ema2 acts as a SUMO E3 ligase}

Ema2 contains an SP-RING domain, although with atypical features: the first cysteine of the zinc ion-binding residues is replaced by histidine, and some stabilizer residues may be missing compared to canonical SP-RING domains \citep{duan2009,yunus2009} (Figure~4A). We found that Ema2 directly interacts with the \textit{Tetrahymena} SUMO E2 enzyme Ubc9 in vitro (Figure~4B), confirming that Ema2 is a bona fide SUMO E3 ligase.

To examine the role of Ema2 in SUMOylation, we expressed HA-tagged Smt3 in an \textit{EMA2} KO strain. HA-Smt3 could replace the essential function of endogenous Smt3 (Figure~EV1). Western blotting using an anti-HA antibody showed that high molecular weight SUMOylated proteins (mainly $>$200~kDa) detected in the WT cross at 4.5 and 6~hpm were reduced to approximately 50\% in the \textit{EMA2} KO cross (Figure~4C). These results indicate that Ema2 is the major SUMO E3 ligase during mid-conjugation stages.

\subsection{Ema2 is required for SUMOylation of Spt6}

To identify SUMOylation targets of Ema2, we used nickel-NTA purification of His-tagged Smt3-conjugated proteins followed by mass spectrometry. While most proteins were detected similarly between the WT cross and \textit{EMA2} KO cross, Spt6, the most abundantly detected protein in the WT cross, was detected at a much lower level in the \textit{EMA2} KO cross (Figure~5A), suggesting that Spt6 is the major Ema2 target for SUMOylation.

Confirmation came from western blotting of HA-tagged Spt6: slower migrating species of Spt6-HA, representing SUMOylated forms, were detected in the WT cross but barely detectable in the \textit{EMA2} KO cross at both 4.5 and 6~hpm (Figure~5B). Immunoprecipitation followed by anti-Smt3 western blotting confirmed that the slower migrating species were SUMOylated Spt6 (Figure~5C).

\subsection{Ema2 is required for the accumulation of lncRNAs in the parental MAC}

Spt6 is a conserved transcription regulator known to promote transcription elongation. We hypothesized that Ema2-dependent Spt6 SUMOylation promotes pMAC-lncRNA transcription. RT-PCR analysis showed that pMAC-lncRNAs were detected in wild-type cells but not in \textit{EMA2} KO cells at 6~hpm for three tested loci (Figure~6A).

Immunostaining with the long double-stranded RNA-specific J2 antibody \citep{schoenborn1991} showed that in wild-type cells, long dsRNAs were detected first in the MIC at 3~hpm, then in the parental MAC at 6~hpm, and in the new MAC at 9~hpm (Figure~6B). In \textit{EMA2} KO cells, long dsRNAs were detected in the MIC at 3~hpm and in the new MAC at 9~hpm but were undetectable in the parental MAC at 6~hpm, indicating that Ema2 is specifically required for pMAC-lncRNA accumulation.

Consistent with the role of pMAC-lncRNAs in recruiting PRC2, H3K27me3 was undetectable in the parental MAC in the \textit{EMA2} KO cells at 6~hpm (Figure~6C), as in the \textit{EMA1}, \textit{TWI1}, and \textit{EZL1} KO cells.

\subsection{Ema2 facilitates chromatin association of Spt6 and RNA polymerase II}

Immunofluorescence staining showed that the MAC localization of both Spt6-HA and Rpb3 (the third largest subunit of RNAPII) was not affected in the absence of Ema2 at the cytological level (Figure~7A). However, cell fractionation experiments revealed a critical difference at the subnuclear level. In the WT cross, Spt6-HA was detected in both the soluble (S1) and chromatin-bound (S2) fractions, with SUMOylated Spt6-HA exclusively in the chromatin-bound fraction. RNAPII showed a similar distribution. In the \textit{EMA2} KO cross, both Spt6-HA and Rpb3 in the chromatin-bound fraction were greatly reduced (Figure~7B). This chromatin association of Spt6 and RNAPII does not require RNA, as they were retained after RNase~A treatment (Figure~7C). These findings were validated in replicated experiments (Figures~EV5, EV6).

\subsection{Spt6 SUMOylation is not necessary for lncRNA transcription}

To directly examine the importance of Spt6 SUMOylation, we generated K-to-R mutants. Among the constructs tested, the HA-SPT6-C-KR mutant showed no detectable SUMOylation (Figure~8B). Surprisingly, cell fractionation showed that Spt6 and RNAPII were detected in the chromatin-bound fraction in HA-SPT6-C-KR cells as in HA-SPT6-WT cells (Figure~8C), and long dsRNAs were detected in the parental MAC of C-KR cells (Figure~8D). These results suggest that Spt6 SUMOylation per se is not required for pMAC-lncRNA transcription. Nonetheless, exconjugants from the HA-SPT6-C-KR cells showed a mild DNA elimination defect (Figure~8E), indicating that Ema2-directed Spt6 SUMOylation plays a role in DNA elimination beyond promoting lncRNA transcription.

\section{Discussion}

In this study, we have identified the SUMO E3 ligase Ema2 as a factor that promotes lncRNA transcription from the parental MAC during conjugation in \textit{Tetrahymena}. The accumulation of these lncRNAs is a prerequisite for target-directed small RNA degradation, which in turn is essential for the selective elimination of germline-limited DNA sequences from the developing somatic genome.

Our findings reveal that Ema2 promotes the chromatin association of both Spt6 and RNAPII, and this activity is essential for pMAC-lncRNA production. The discovery that Spt6 is the major SUMOylation target of Ema2 suggests a model in which SUMOylation modulates chromatin-associated transcription. However, the unexpected finding that a SUMOylation-defective Spt6 mutant (C-KR) retains the ability to support lncRNA transcription indicates that the relationship is more complex: Ema2 likely promotes lncRNA transcription through SUMOylation of targets other than Spt6, or through non-catalytic mechanisms.

The involvement of a SUMO pathway in DNA elimination appears to be conserved among ciliates, as RNAi knockdown of the SUMO E1 enzyme UBA2 and SUMO in \textit{Paramecium} also impairs DNA elimination \citep{matsuda2006}. Furthermore, SUMOylation has recently been linked to the piRNA pathway in \textit{Drosophila} and \textit{C.\ elegans}, where the SUMOylation of Panoramix connects piRNA signaling to heterochromatin establishment \citep{zhao2022}. Thus, the regulatory role of SUMOylation in small RNA-directed chromatin silencing may be a broadly conserved feature across eukaryotes.

In contrast to the passive mechanism of lncRNA production in fission yeast, where eviction of the HP1 protein Swi6 during mitosis allows temporary transcription \citep{chen2008,kloc2008}, and the active mechanism in \textit{Drosophila}, where Rhino recruits dedicated transcription factors \citep{andersen2017}, \textit{Tetrahymena} employs a distinct mechanism through SUMOylation-dependent promotion of transcription factor chromatin association. This represents a third paradigm for how organisms ensure transcription of heterochromatic regions that are otherwise targeted for silencing.

\section{Materials and Methods}

\subsection{Strains and cell culture}
Wild-type \textit{Tetrahymena thermophila} strains B2086, CU427, and CU428, and the MIC-defective strains B*VI and B*VII were obtained from the Tetrahymena Stock Center. Cells were grown in SPP medium \citep{gorovsky1975} containing 2\% proteose peptone at 30$^\circ$C. For conjugation, growing cells of two different mating types were washed, pre-starved, and mixed in 10~mM Tris (pH~7.5) at 30$^\circ$C.

\subsection{Molecular tagging and knockout strains}
Epitope-tagged strains (EMA2-HA, SPT6-HA, HA-SMT3, His-SMT3) and knockout strains were generated using established biolistic transformation methods. Constructs expressing K-to-R substitution mutants of Spt6 (DOL-KR, N-KR, M-KR, C-KR) were introduced into SPT6 KO cells for genetic rescue experiments.

\subsection{Immunofluorescence staining and microscopy}
Cells were fixed with paraformaldehyde and processed for immunofluorescence staining with the indicated antibodies. DNA was stained with DAPI. For dsRNA detection, the J2 antibody specific for long double-stranded RNA was used \citep{schoenborn1991}.

\subsection{Western blotting and immunoprecipitation}
Total proteins were harvested at the indicated time points and analyzed by western blotting with the appropriate antibodies. For SUMOylated protein detection, HA-tagged Smt3-conjugated proteins were analyzed with anti-HA antibodies. For mass spectrometry analysis of SUMOylation targets, His-Smt3-conjugated proteins were purified using Ni-NTA beads.

\subsection{Cell fractionation}
Conjugating cells were incubated with lysis buffer to release cytoplasmic and nucleoplasmic proteins (S1 fraction). The insoluble fraction was sonicated to release chromatin-bound proteins (S2 fraction). For RNase treatment experiments, the insoluble fraction was incubated with or without RNase~A before fractionation.

\subsection{Small RNA analysis}
Total RNA was isolated from conjugating cells and analyzed by northern blotting using the radioactive Mi-9 probe. For genome-wide analysis, small RNAs (26--32~nt) were sequenced and mapped to genomic tiles of MDSs and IESs.

\subsection{RT-PCR for lncRNA detection}
Total RNA from conjugating cells at 6~hpm was used for RT-PCR with primers designed to amplify lncRNAs from the MAC M, L8, and R2 loci. RPL21 mRNA was used as a positive control.

\subsection{DNA-FISH}
Exconjugants at 36~hpm were analyzed by DNA-FISH using fluorescent probes complementary to the Tlr1 transposable element.

\subsection{In vitro binding assay}
GST-tagged Ema2 and His-tagged Ubc9 were recombinantly expressed in \textit{E.\ coli}. GST-Ema2 was immobilized on glutathione beads and incubated with His-Ubc9. Bound proteins were analyzed by western blotting with anti-GST and anti-His antibodies.

\section*{Acknowledgments}
We thank Dr.\ James Forney (Purdue University) for the anti-Smt3 antibody. We acknowledge the Tetrahymena Stock Center for providing strains.

\bibliography{references}

\end{document}
